\section{TSMC 与半导体产业}

\subsection{TSMC 的垄断地位}

TSMC 生产世界上大部分先进芯片。代工模式（foundry model）的成功依赖于规模经济——当你为所有客户制造芯片时，每个制程节点的研发成本被摊薄到整个行业。

由张忠谋创立——他在德州仪器（TI）被跳过 CEO 后转战台湾，创造了改变世界的商业模式。这是一个深刻的历史偶然：如果张忠谋当上了 TI CEO，台湾可能不会成为全球半导体中心。

\subsection{台湾的独特优势}

台湾半导体产业的文化基因：

\begin{itemize}
    \item \textbf{人才密度}：顶尖毕业生以 \$70-80K 年薪加入 TSMC——在美国同等水平的工程师薪资是 3-5 倍
    \item \textbf{工作伦理}：极端的投入度，远超硅谷标准
    \item \textbf{应急响应}：地震后的快速恢复能力，整个产业链的抗风险弹性
    \item \textbf{制造纪律}：半导体制造需要持续数月的精确控制，台湾文化适合这种工作
\end{itemize}

全球只有三个地方在做尖端研发：新竹（TSMC）、Hillsboro, Oregon（Intel）、Pyeongtaek（Samsung）。这个名单可能还会缩短——Intel 和 Samsung 都在挣扎。

\begin{figure}[H]
\centering
\includegraphics[width=0.92\textwidth]{frame-02-tsmc.jpg}
\caption{访谈进入 TSMC 与地缘风险讨论\protect\footnotemark}
\end{figure}
\footnotetext{视频画面时间区间：01:31:05--01:31:16。}

\begin{warningbox}{台湾：全球科技的最大单点故障}
\textit{``如果我有几枚导弹，我确切知道哪里能造成最大的经济损害——就是 TSMC、Intel、Samsung 的研发中心。''} (Dylan)

美国 CHIPS Act 只有 \$50B——对比中国每年约 \$200B 的半导体补贴。而 Intel 正在衰落：失去制程领先地位、没有 AI 芯片胜利、CEO 被解雇、PC 和服务器市场份额流失。Samsung 的制造良率也在下降。

最极端的情景：世界可能最终只剩 TSMC 一家做尖端研发。如果台海发生冲突，全球科技产业将遭受无法估量的损失。
\end{warningbox}

\subsection{本章小结}

半导体供应链高度集中于 TSMC，台湾的地缘政治风险是全球科技的最大单点故障。Intel 的衰落和 Samsung 的困境意味着集中度在增加而非减少。

